% Image d'un objet à l'infini par une lentille convergente : un faisceau incliné converge dans le plan focal image (L'essentiel, chapitre 24)
\begin{tikzpicture}[x=1cm,y=1cm]
\draw[otMuted,line width=0.4pt] (-2.6,0) -- (5.6,0);
\draw[otNavy,line width=1pt,{Stealth[length=2mm]}-{Stealth[length=2mm]}] (0,-1.5) -- (0,1.5) node[pcLegende,anchor=south,text=otNavy] {$L$};
\fill[otInk] (0,0) circle (0.04) node[pcLegende,anchor=north east] {$O$};
\fill[otInk] (4,0) circle (0.04) node[pcLegende,anchor=north] {$F'$};
\fill[otInk] (-4,0) circle (0.0);
\draw[otMuted,line width=0.4pt,dash pattern=on 2pt off 1.5pt] (4,-1.4) -- (4,1.2);
% faisceau incliné : pente -0,2 ; image B' en (4 ; -0,8)
\foreach \y in {1.2,0,-1.2}{
  \draw[otRed!80!black,line width=0.7pt] (-2.4,{\y+0.48}) -- (0,\y) -- (4,-0.8);
  \draw[otRed!80!black,line width=0.7pt,-{Stealth[length=1.5mm]}] (-2.4,{\y+0.48}) -- (-1.2,{\y+0.24});}
\fill[otRed!80!black] (4,-0.8) circle (0.05) node[pcLegende,anchor=west,text=otRed!80!black] {$B'$};
\draw[otMuted,line width=0.4pt] (1.2,0) arc[start angle=0,end angle=-11.3,radius=1.2];
\node[pcLegende,anchor=west] at (1.25,-0.13) {$\theta$};
\draw[otGreenDark,line width=0.5pt,{Stealth[length=1.2mm]}-{Stealth[length=1.2mm]}] (4.45,0) -- (4.45,-0.8) node[pos=0.5,pcLegende,anchor=west,text=otGreenDark] {$f'\tan\theta$};
\end{tikzpicture}
\par\smallskip\legende{Les rayons issus d'un point très éloigné arrivent parallèles entre eux ; ils convergent en un point $B'$ du plan focal image, situé sur le rayon qui passe par le centre optique.}
